\subsection{何尊与大盂鼎：它们说了什么，没有说什么}

\begin{event}{成王时期说（绝对年与王年仍有讨论）}{何尊铸铭与“宅兹中国”}{器铭可靠；王年、释文细节与所指空间有争议}{成周及王室礼仪活动的语境}\eventanchor{WZ-HE-ZUN}{西周·何尊铸铭}
一件何尊使成周的政治语言有了可以触摸的载体。器铭追记周王“初迁宅于成周”，并把王的命令、臣下的服务与“余其宅兹中国”放在同一段记忆里。对受命者家族来说，铸器是把一次进入王室秩序的经历交给子孙；对今天的读者来说，它让\eventref{WZ-CHENG-ROYAL-CEREMONY}{西周·成周告成与册命}不必完全依靠后世叙事。

然而，何尊不是一张城图，也不是周初封国总表。“中国”在这里指向王权所构想的中心空间，不能直接换算为现代国土；铭文也没有交代城墙长度、迁置人数或某一封国的边界。王究竟应系成王还是其他早期周王、器文中个别字句如何断读，以及“初迁”处于经营东方的哪一阶段，均仍须由金文断代、传世文献和洛阳考古相互检验。器物的分量正在于它能收紧我们的判断，而不是许可我们把材料未说出的细节补满。
\end{event}

何尊常被视为理解成周的关键，因为铭文记述“王初迁宅于成周”，并有“余其宅兹中国”的著名语句。它让我们看见，一位周王及其近臣确实把成周与“中”这一政治空间联系起来；这里的“中国”是王权所设想的中心区域，不是现代民族国家的名称。与主要经过后世整理的叙事相比，这是一件接近西周早期政治语境的材料。

然而，何尊也不能独自画出成周蓝图。铭文的释读、器物的年代和所述“初迁”在王室经营东方过程中的位置，都有研究空间；它没有列出城的边界、迁入者人数或全部封国。把它当作周初建都的实物证据是恰当的，把它扩展成一份完整的都市规划书则超出了证据。

大盂鼎同样重要，却与成周问题的距离更远。它记载王对盂的训诰，并以殷亡的教训警戒受命者；这使它成为观察周初或早期西周朝廷如何用商亡来解释自我约束、官员责任和受命秩序的珍贵材料。它不是一份成周工程记录，也不能单凭鼎文证明某次迁置或某一封国的设立。将大盂鼎直接用作“周初一切制度细节”的证据，正会抹掉它的文类和用途。

两器的价值在于互相校正材料层次：何尊把成周带入同时代的政治表达，大盂鼎显示王室如何把统治得失写进贵族训诰；\sourcebook{尚书}诸篇和后出的史书则提供更连贯、也更需要辨析成书背景的叙事。金文不是没有立场的档案，而是为纪念、传承和展示而铸造的文字。它让我们更接近王室与贵族的行动语言，却不直接替我们说出所有被征调、迁置和服役者的经验。

何尊的形制和文字还提示了一层制度条件：王命、祭祀与赏赐要成为可继承的政治记忆，须有贵族愿意并有能力铸造、保存和展示礼器。器物于是把一次受命延长为家族的祖先叙事。它能说明成周在王室近臣的行动语言中占有位置，却不能代表未铸器者，更不能把一位受命者的纪念推广为整个东方社会的共同感受。

因此，金文的价值不在于替代其他材料，而在于迫使叙述保持比例。与\sourcebook{尚书·召诰}、\sourcebook{洛诰}的建城和训诫语言并读时，何尊支持早周王室确曾以成周组织政治空间；与洛阳考古并读时，它又提醒人们城市、礼仪和王命彼此依赖。至于“初迁”究竟是哪一次行动、相关王年如何换算、所谓“中”涵盖多大空间，仍应让释读、断代与城市考古的争论留在判断之内，而非以一件名器关闭问题。

\begin{sourcenote}[以不同材料互证，而不让一件器物承担全部证明]
何尊、大盂鼎的释文、断代和历史语境一直是金文研究的重要议题；《尚书》周初篇章的文本形成和可靠性也有长期争论。这里采用的不是“金文必真、传世文献必伪”的简单排序，而是分别追问：它在何种场合被制作或编纂？它直接证明什么？哪一步需要由其他材料补足？这种读法能支持成周与东方经营的基本轮廓，也要求我们对具体年份、人数、城址范围和动机保留判断。
\end{sourcenote}

\begin{turningpoint}[制度得失：东方网络的收益、账单与后果]
\term{它解决了什么}：成周缩短了王室接近东方的距离，迁置与封国重组把征服后的不确定关系纳入可管理的秩序。以地方精英承担日常治理，提升了王室在广阔空间中的动员、礼仪整合和军事反应能力。

\term{谁承担成本}：筑城、运输、驻守、祭祀和迁置所需的人力与物资，最终要由王室、封国及其辖下人群共同供给。材料不能为这些负担开出精确账单，却足以表明政治整合并非无代价；尤其被重新安置的人群，很难在现有材料中留下自己的陈述。

\term{它留下什么副作用}：王室以承认地方精英来换取治理能力，也帮助他们积累土地、人口、武力和代际声望。当王室仍能册命、赏赐并调解冲突时，这是一种制度收益；当这些能力下降时，原先的协作渠道会转为各方自行结盟的资源。到\eventref{WZ-0771-HAOJING}{西周·镐京陷落}后的东迁，洛邑仍可承载周王室的礼仪名分，却已不能恢复西周初年那种以资源和军力支撑的网络。
\end{turningpoint}
